\section{我的学习认识与总结}
\subsection{从记住工具名称到理解技术边界}
回顾这次任务，我的认识主要发生在三个方面。首先，前端技术之间有层次：HTML、CSS 和 JavaScript
是基础，Vue 帮助组织界面，TypeScript 与 Vite 改善工程过程，Three.js 和编辑器则服务于具体场景。%
它们不是越多越好，而是需要各自承担清楚的职责。

其次，相似的页面表现可能来自不同原理。卡片倾斜并不一定需要实体模型，%
懒加载代码并不等于图片或数据也已优化，能输入富文本也不代表已经拥有结构化文档或协同编辑。%
遇到需求时先确定要解决什么问题，再选择恰当层次的技术，比直接套用听起来先进的方案更有效。

最后，展示效果与稳定性需要一起考虑。三维效果要尊重性能和减少动画偏好；内容加载要处理失败与旧响应；%
编辑能力要保护文档结构和已有修改；图片和下载还受到安全与权限边界约束。%
对初学者来说，这些约束起初可能显得繁琐，实际却使前端从一次演示变成可以继续维护的应用。

\subsection{结合本项目的后续学习方向}
我的后续学习会继续沿着已有任务展开。三维方面，先深入理解坐标、透视和帧时间，%
再在确有模型展示需求时研究 glTF、WebGL 或 WebGPU；内容加载方面，先完善资源规格和测量，%
再评估公开文章的预渲染；编辑方面，优先研究输入法、撤销与文档往返，再判断是否迁移到结构化编辑器。%
这样的学习顺序能够把新知识与已有问题对应起来。

我也需要学习如何验证自己的判断。类型检查可以发现一类代码问题，组件测试可以检查状态变化，%
真实浏览器才能更直接地观察布局、焦点和输入体验。项目既有验证记录与本次技术调研承担不同作用；%
我没有把阅读官方示例写成实际部署，也没有把预览截图写成真实业务已经全部验收。%
对技术优缺点的认识，应在后续实践中继续被检验。

这份报告从前端入门出发，以 WEMOVE 的具体任务连接了三维展示、内容加载和内容编辑的现状。%
对我最有价值的收获，是能够解释一项技术为什么适合某个任务、又在哪些地方需要付出代价。%
前端的发展提供了更丰富的工具，而学习的目标是让这些工具帮助用户更顺畅地阅读、操作和维护内容。
